% Generated from docs/REPORT_DRAFT.md; edit the manuscript then rerun the converter.
\chapter{Công trình và nền tảng liên quan}

\section{Phương pháp khảo sát}

Khảo sát tập trung vào ba cụm: Knowledge Graph construction/entity resolution,
KG question answering và explainable KG recommendation. Từ khóa gồm “knowledge
graph question answering”, “knowledge graph recommender systems”, “explainable
graph recommendation”, “movie knowledge graph recommendation” và “temporal
knowledge graph QA”. Nguồn ưu tiên là W3C, tài liệu chính thức của nền tảng,
proceedings hội nghị, tạp chí peer-reviewed và bản paper của tác giả. Khoảng thời
gian ưu tiên là 2023–2026; nguồn nền tảng cũ hơn chỉ dùng cho định nghĩa chuẩn.

Tiêu chí đưa vào gồm: liên quan trực tiếp tới KGQA/recommendation/construction;
có mô tả phương pháp hoặc evaluation; có metadata xuất bản xác minh được. Bài
blog, trang tổng hợp và nội dung không có phương pháp rõ ràng bị loại khỏi survey
cốt lõi.

\section{Tổng hợp nghiên cứu}

Guo và cộng sự khảo sát các hệ gợi ý dựa trên Knowledge Graph, phân loại cách
khai thác graph và nhấn mạnh hai mục tiêu accuracy và explainability \cite{guo2023kgsurvey}. Kết quả
này hỗ trợ quyết định đánh giá đồng thời ranking metric và explanation coverage.
Các nghiên cứu về explainable graph recommendation cho thấy đường/feature graph
có thể cung cấp lý do dễ kiểm tra hơn latent factor \cite{caro2023graph}, \cite{colas2023knowledge}. Tuy nhiên, nhiều
phương pháp tối ưu personalization từ user–item interaction hoặc dùng embedding;
đề tài này không có user history nên chỉ tuyên bố item-to-item recommendation.

Các công trình conversational recommendation dùng explicit reasoning chain để
tạo recommendation rationale \cite{ren2024explicit}. Đề tài chia QA và recommendation thành hai
endpoint, nhưng chia sẻ nguyên tắc: mọi kết quả phải có graph evidence. Với miền
phim, các nghiên cứu kết hợp KG với time series hoặc topic model nhằm tăng độ
chính xác \cite{zhang2023movie}, \cite{saat2024enhanced}. Đây là baseline khái niệm quan trọng, đồng thời chỉ ra giới
hạn của đề tài: chưa mô hình hóa thời gian tương tác hoặc nội dung overview bằng
topic model.

Khảo sát Temporal KGQA năm 2024 phân biệt câu hỏi trên fact thay đổi theo thời
gian với graph tĩnh \cite{su2024temporal}. Movie KG hiện chỉ lưu \texttt{release\_date} như datatype
property để filter/sort; nó chưa phải Temporal KG và không hỗ trợ valid time hay
transaction time. BYOKG dùng LLM-backed symbolic agent để khám phá KG chưa thấy,
sinh query-program exemplar và tổng hợp chương trình truy vấn zero-shot \cite{agarwal2024byokg}.
Khác với hướng đó, đề tài cố định schema/operation whitelist để ưu tiên safety và
reproducibility trong phạm vi Movie KG. Vì vậy smoke test 10 câu của đề tài chỉ
được dùng để kiểm tra hệ thống, không gọi là benchmark production.

\section{Khoảng trống và vị trí đề tài}

Đề tài không cố cạnh tranh với KGQA/recommender quy mô nghiên cứu. Khoảng trống
được chọn là một workflow học phần nhưng đầu cuối và kiểm chứng được: thu thập
đa nguồn, identity/provenance, hai biểu diễn graph, suy diễn có evidence, API/UI,
test và artifact thực nghiệm. Điểm khác biệt quan trọng là LLM không trực tiếp
sinh Cypher hoặc câu trả lời, recommendation score được giải thích bằng các quan
hệ trong graph, và các kết quả đo đều gắn với snapshot cùng cấu hình cụ thể.

\section{So sánh có cấu trúc các công trình liên quan}

\begin{longtable}{|>{\raggedright\arraybackslash}p{0.18\textwidth}|>{\raggedright\arraybackslash}p{0.18\textwidth}|>{\raggedright\arraybackslash}p{0.18\textwidth}|>{\raggedright\arraybackslash}p{0.18\textwidth}|>{\raggedright\arraybackslash}p{0.18\textwidth}|}
\caption{Tổng hợp so sánh có cấu trúc các công trình liên quan}\label{tab:3-1} \\
\hline
\textbf{Công trình} & \textbf{Dữ liệu/tín hiệu} & \textbf{Phương pháp} & \textbf{Explanation} & \textbf{Khác biệt với đề tài} \\ \hline
\endfirsthead\hline \textbf{Công trình} & \textbf{Dữ liệu/tín hiệu} & \textbf{Phương pháp} & \textbf{Explanation} & \textbf{Khác biệt với đề tài} \\ \hline\endhead
Guo et al. \cite{guo2023kgsurvey} & nhiều benchmark & survey taxonomy & phân tích nhiều họ & dùng để định vị, không phải baseline chạy \\ \hline
Caro-Martínez et al. \cite{caro2023graph} & interaction graph & link prediction/common neighbors & graph-interpretable & đề tài dùng metadata graph, không user interaction \\ \hline
Colas et al. \cite{colas2023knowledge} & user–item + item KG & recommendation/NLG & natural-language grounding & đề tài trả evidence xác định, không NLG \\ \hline
Ren et al. \cite{ren2024explicit} & conversational recommendation & KG reasoning chain & explicit chain & QA stateless, recommendation tách endpoint \\ \hline
Zhang et al. \cite{zhang2023movie} & rating + KG + time series & collaborative filtering & quan hệ KG & không có temporal user signal \\ \hline
Saat et al. \cite{saat2024enhanced} & MovieLens + topic/metadata & content-based KG & feature/profile & dùng TMDB, không topic model \\ \hline
Su et al. \cite{su2024temporal} & temporal KGQA studies & survey taxonomy & temporal reasoning & release date mới là property \\ \hline
Agarwal et al. \cite{agarwal2024byokg} & unseen/domain KG & LLM program synthesis & query program/path & fixed schema/whitelist compiler \\ \hline
\end{longtable}

“Dùng Knowledge Graph” không phải một phương pháp duy nhất. KG có thể là nguồn
side information cho model, không gian reasoning, cấu trúc explanation hoặc
operational database. Trong đề tài, graph đồng thời là operational store và
evidence substrate; model học sâu không trực tiếp quyết định ranking hoặc fact.

\section{Liên hệ survey với quyết định thiết kế}

Survey dẫn đến bốn quyết định. Thứ nhất, explainability phải được định nghĩa bằng
evidence có thể kiểm tra, không chỉ text nghe hợp lý. Thứ hai, metric recommendation
phải đi cùng corpus/rubric và K. Thứ ba, QA sinh query tự do có flexibility nhưng
execution risk cao; constrained plan phù hợp phạm vi học phần hơn. Thứ tư, các
phương pháp personalized không được dùng để quảng bá hệ thống không có user profile.
